% Ad Familiares 8.3 (ad-familiares-08-03)
% Source: https://raw.githubusercontent.com/PerseusDL/canonical-latinLit/master/data/phi0474/phi056/phi0474.phi056.perseus-lat1.xml
% Fetched by scripts/fetch_latin.py

% Dateline (Perseus): Scr. Romae circ. iv Id. Iun. a. 703 (51).

\ciceroLetterOpener{CAELIVS CICERONI S.}

\ciceroSection{1} estne ? vici et tibi saepe, quod negaras discedens curaturum tibi, litteras mitto? est , si quidem perferuntur, quas do. atque hoc eo diligentius facio quod, cum otiosus sum, plane ubi delectem otiolum meum non habeo. tu cum Romae eras, hoc mihi certum ac iucundissimum vacanti negotium erat, tecum id oti tempus consumere idque non mediocriter desidero, ut mihi non modo solus esse sed Romae te profecto solitudo videatur facta, et, qui, quae mea neglegentia est, multos saepe dies ad te, cum hic eras, non accedebam, nunc cotidie non esse te ad quem cursitem discrucior. maxime vero ut te dies noctesque quaeram competitor Hirrus curat. quo modo illum putas auguratus tuum competitorem dolere et dissimulare me certiorem quam se candidatum? de quo ut quem optas quam primum nuntium accipias, tua medius fidius magis quam mea causa cupio; nam mea, si fio, forsitan cum locupletiore referam; sed hoc usque eo suave est ut, si acciderit, tota vita risus nobis desse non possit. † sed tanti sed me hercules non multum M. Octavius eorum odia quae Hirrum premunt, quae permulta sunt, sublevat.

\ciceroSection{2} quod ad Philotimi liberti officium et bona Milonis attinet, dedimus operam ut et Philotimus quam honestissime Miloni absenti eiusque necessariis satis faceret et secundum eius fidem et sedulitatem existimatio tua conservaretur.

\ciceroSection{3} illud nunc a te peto, si eris, ut spero, otiosus, aliquod ad nos, ut intellegamus nos tibi curae esse, su/ntagma conscribas. ' qui tibi istuc,' inquis, 'in mentem venit, homini non inepto?' opto aliquod ex tam multis tuis monumentis exstare quod nostrae amicitiae memoriam posteris quoque prodat. cuius modi velim, puto, quaeris. tu citius, qui omnem nosti disciplinam, quod maxime convenit excogitabis, genere tamen quod et ad nos pertineat et didaskali/an quandam, ut versetur inter manus, habeat.
